\documentclass[10pt]{article}
\usepackage[a4paper,margin=0.6in]{geometry}
\usepackage{amsmath,amssymb,amsfonts,mathtools}
\setlength{\jot}{4pt}
\usepackage{xcolor}
\setlength{\parindent}{0pt}
\setlength{\parskip}{5pt}
\newcommand{\Leps}{L_{\epsilon}}
\newcommand{\KK}{\mathcal K(x-y)\!\cdot\!\mathcal K(x'-y')}
\newcommand{\Wbr}{\left[U^{ce'}(y') - U^{ce'}(x')\right] \left[U^{ce}(y) - U^{ce}(x)\right] \rho^{e'}(x')\, \rho^{e}(x)}
\newcommand{\Wz}{\left[U^{ce'}(z) - U^{ce'}(x')\right]\left[U^{ce}(z) - U^{ce}(x)\right]\rho^{e'}(x')\,\rho^{e}(x)}
\title{Group II, Row IV: extraction of the BFKL and finite terms\\ from the left-over expression}
\author{}
\date{}
\begin{document}
\maketitle
\vspace{-2.5em}

\section{The left-over expression and its check}
After the JIMWLK term and the DGLAP $+$ RC terms have been extracted, the left-over is
\begin{align}\label{eq:leftover}
\frac{d^{3}N_{\rm left}}{d^{3}k}  & =
-\,\frac{1}{(2\pi)^{3}}\,\frac{g^{4}\,N_{c}}{8\pi^{4}}\,\frac{1}{k^{+}}\,\log\Big(1-\frac{k^{+}}{\vee}\Big)\,\int_{x,y,x',y'}e^{-ik\cdot(y'-y)}\,\frac{(x-y)\cdot(x'-y')}{(x-y)^{2}(x'-y')^{2}}\,\Big[\frac{2}{\epsilon}+2\gamma_{E}+\log\frac{(y'-y)^{2}\mu^{2}_{\overline{MS}}}{4}\Big]\notag \\ & \hspace{2em}\times\Wbr\notag\\
&+\frac{g^{2}}{4\pi^{2}}\frac{1}{(2\pi)^{3}}\frac{g^{2}}{2\pi^{2}}\frac{N_{c}}{k^{+}}\int_{k^{+}/\vee}^{1}d\xi\,\frac{1-\xi}{\xi}\int_{x,y,x',y'}e^{-ik\cdot(y'-y)\frac{1}{\xi}}\Bigg[\frac{(x-y)\cdot(x'-y')}{(x-y)^{2}(x'-y')^{2}}-2\,\frac{(x-y)\cdot(y'-y)\,(x'-y')\cdot(y'-y)}{(x-y)^{2}(x'-y')^{2}(y'-y)^{2}}\Bigg]\notag \\ & \hspace{2em}\times\Wbr\,.
\end{align}
Write $C\equiv\frac{1}{(2\pi)^{3}}\frac{g^{4}N_{c}}{8\pi^{4}}\frac{1}{k^{+}}$ (both prefactors
of \eqref{eq:leftover} equal $C$, since
$\frac{g^{2}}{4\pi^{2}}\cdot\frac{g^{2}}{2\pi^{2}}N_{c}=\frac{g^{4}N_{c}}{8\pi^{4}}$),
$\Leps=\frac{2}{\epsilon}+2\gamma_{E}+\log\frac{(y'-y)^{2}\mu^{2}_{\overline{MS}}}{4}$, and
$\bar\xi=\frac{k^{+}}{\vee}$.

\textbf{The left-over is correct}, and its two terms are identified as follows.
\begin{itemize}
\item[(a)] The first term is the \emph{endpoint piece} $-C\,h(1)\,\log(1-\bar\xi)$ of the soft
channel, $h(1)=\big[\frac{d_{\perp}}{2}\Leps+1\big]e^{-ik(y'-y)}=\Leps\,e^{-ik(y'-y)}$: the full
endpoint logarithm was $\log\frac{(\vee-k^{+})k^{+}}{\vee\Lambda}
={\log\frac{k^{+}}{\Lambda}}+\log(1-\bar\xi)$, the $\log\frac{k^{+}}{\Lambda}$ part left with
the JIMWLK term, and $\log(1-\bar\xi)$ stayed. (Checked numerically: soft channel minus JIMWLK
minus the locally subtracted integral equals $-C\,h(1)\log(1-\bar\xi)$ to $O(\Lambda)$.)

\emph{Convention remark.} This term is present because the DGLAP term was extracted in the
\emph{locally subtracted} plus-prescription convention,
$\int_{\bar\xi}^{1}d\xi\,\frac{h(\xi)-h(1)}{1-\xi}$. If instead the \emph{standard} plus
distribution on $[0,1]$ is used in $\mathcal P_{gg}$ --- the convention in which
$\mathcal P_{gg,+}=2N_{c}\big[\xi(1-\xi)+\frac{1-\xi}{\xi}+\frac{\xi}{(1-\xi)_{+}}\big]$ is the
textbook splitting function --- then
$\int_{\bar\xi}^{1}\frac{h}{(1-\xi)_{+}}=\int_{\bar\xi}^{1}\frac{h-h(1)}{1-\xi}
+h(1)\log(1-\bar\xi)$ absorbs this first term into the DGLAP integral, and it disappears from
the finite block. The two bookkeepings differ by exactly this one term; nothing else changes.
\item[(b)] The second term is the \emph{azimuthal} term: the $\Leps$-free part of the regrouped
expression, $+C\int_{\bar\xi}^{1}d\xi\,\frac{1-\xi}{\xi}
\big[\KK-d_{\perp}\frac{(\mathcal K\cdot r)(\mathcal K'\cdot r)}{r^{2}}\big]f(\xi)$,
$r=y'-y$, with $d_{\perp}\to2$ (legitimate, since no pole multiplies it) and the upper limit
already extended to $1$ at a cost $O(\Lambda)$.
\end{itemize}

\section{Step 1: no rapidity logarithm hides in the left-over}
Before anything can be ``extracted'', note that \emph{neither term of \eqref{eq:leftover}
contains a $\log\frac{\vee}{k^{+}}$}:
\begin{itemize}
\item[(i)] The coefficient of the first term is $\log(1-\frac{k^{+}}{\vee})$, which tends to
$0$ as $\vee\to\infty$ at fixed $k^{+}$. It does not grow with the rapidity interval; it is a
finite endpoint (convention) term, not a rapidity logarithm.
\item[(ii)] The second term is finite as $\bar\xi\to0$. Two mechanisms protect it.
First, at fixed positions, substituting $u=\frac1\xi$ gives
\begin{align}
\int_{\bar\xi}^{1}d\xi\,\frac{1-\xi}{\xi}\,e^{-ik\cdot r/\xi}
=\int_{1}^{1/\bar\xi}du\,\Big(\frac{1}{u}-\frac{1}{u^{2}}\Big)\,e^{-ik\cdot r\,u}\,,
\end{align}
which \emph{converges} as $\frac{1}{\bar\xi}=\frac{\vee}{k^{+}}\to\infty$ for $r\neq0$: the
phase oscillates and the integral approaches a finite limit (checked numerically: the value
stabilizes while $\log\frac{\vee}{k^{+}}$ would have grown without bound). Second, the
dangerous collapse region $r=O(\xi)$, where the phase stops oscillating, is doubly suppressed:
writing $y'=y+\xi s$ costs a factor $\xi^{2}$ from the measure, and the azimuthal bracket has
\emph{exactly zero} angular average over the direction $\hat r$,
\begin{align}
\int_{0}^{2\pi}\frac{d\theta_{\hat r}}{2\pi}
\Big[\,\mathcal K\!\cdot\!\mathcal K'-2\,(\mathcal K\!\cdot\!\hat r)(\mathcal K'\!\cdot\!\hat r)\Big]
=\mathcal K\!\cdot\!\mathcal K'-\mathcal K\!\cdot\!\mathcal K'=0
\end{align}
(in two dimensions $\langle2\hat r^{i}\hat r^{j}\rangle_{\theta}=\delta^{ij}$).
\end{itemize}
\textbf{Conclusion of Step 1.} The BFKL logarithm cannot be found \emph{inside}
\eqref{eq:leftover}. It lives in the already-extracted DGLAP term: the soft channel
$2N_{c}\frac{1-\xi}{\xi}\otimes\Leps$ produces a hidden $\log\frac{\vee}{k^{+}}$ at its
$\xi\to0$ endpoint ($\xi\to0$ is $p^{+}\to\vee$: the unmeasured gluon carries almost all of the
projectile momentum and the measured gluon is the soft one --- the projectile-side rapidity
region), and a rapidity logarithm multiplying $\Leps$ is not physical. The BFKL term is
therefore produced by an \emph{exact add-and-subtract}: a subtraction term with the same
$\xi\to0$ limit is added to the finite block (where it protects the DGLAP soft channel), and
its exactly computed value is subtracted back. Adding zero changes nothing, so the construction
may equally well be carried out on the left-over \eqref{eq:leftover}, which is what we now do.

\section{Step 2: constructing the subtraction term, explicitly}

\textbf{(2a) Two building blocks.} Throughout, write
\begin{align}
W(u,u')\equiv\left[U^{ce'}(u')-U^{ce'}(x')\right]\left[U^{ce}(u)-U^{ce}(x)\right]\rho^{e'}(x')\,\rho^{e}(x)\,,
\end{align}
so that the colour factor of \eqref{eq:leftover} is $W(y,y')$, and define the $\xi$-dependent
LO structure
\begin{align}\label{eq:dnxilo}
\frac{d^{3}N_{\xi\,\rm LO}}{d^{3}k} = \frac{1}{(2\pi)^{3}}\, \frac{g^{2}}{2\pi^{2}}\, \frac{1}{k^{+}}\,\int_{x,y,x',y'}\, e^{-ik\cdot (y'-y)\,\frac{1}{\xi}}\, \KK\;W(y,y')
\end{align}
($\xi=1$ is the LO cross-section itself). The second building block is the dimensionally
regularized transverse integral derived in the companion reduction note,
\begin{align}\label{eq:zint}
\int_{z}^{(d_{\perp}=2-\epsilon)} K(y,y',z)=-\pi\,\Leps\,,\qquad
K(a,b,z)=\frac{(a-z)\cdot(b-z)}{(a-z)^{2}(b-z)^{2}}\,,
\end{align}
which we read \emph{backwards}:
\begin{align}\label{eq:lepsasz}
\Leps=-\frac{1}{\pi}\int_{z}K(y,y',z)\,.
\end{align}
Eq.~\eqref{eq:lepsasz} trades the abstract ultraviolet bracket $\Leps$ for a concrete
transverse integral over one new point $z$. That is the whole trick of Step 2: once $\Leps$ is
a $z$-integral, a counterterm can be built under the \emph{same} integral sign and compared
with it point by point.

\textbf{(2b) The term that needs a partner.} The DGLAP term extracted before this note contains
the soft channel
\begin{align}\label{eq:dsoft}
D_{\rm soft}\;=\;-\,\frac{g^{2}}{8\pi^{2}}\int_{\bar\xi}^{1}\frac{d\xi}{\xi^{2}}\;2N_{c}\,\frac{1-\xi}{\xi}\;\Leps\;\frac{d^{3}N_{\xi\,\rm LO}}{d^{3}k}\,.
\end{align}
It has the two dangerous features at once: the ultraviolet bracket $\Leps$, and a hidden
$\log\frac{\vee}{k^{+}}$ from its $\xi\to0$ endpoint. A cross-section must not contain
$\log\frac{\vee}{k^{+}}\times\Leps$. The subtraction term is built to reproduce \emph{exactly}
the $\xi\to0$ behaviour of \eqref{eq:dsoft} --- nothing more --- so that in the difference the
rapidity logarithm and the $\mu$-dependence disappear together.

\textbf{(2c) Put $D_{\rm soft}$ under a $z$-integral.} Insert \eqref{eq:lepsasz} into
\eqref{eq:dsoft} and write out \eqref{eq:dnxilo}:
\begin{align}\label{eq:dsoftz}
D_{\rm soft}
=+\,\frac{g^{2}}{8\pi^{2}}\int_{\bar\xi}^{1}\frac{d\xi}{\xi^{2}}\,2N_{c}\,\frac{1-\xi}{\xi}\;
\frac{1}{(2\pi)^{3}}\frac{g^{2}}{2\pi^{2}}\frac{1}{k^{+}}
\int_{x,y,x',y',z} e^{-ik\cdot(y'-y)\frac{1}{\xi}}\;\frac{1}{\pi}\,K(y,y',z)\;\KK\;W(y,y')\,.
\end{align}
This is still exact; the overall sign flipped because of the minus sign in \eqref{eq:lepsasz}.

\textbf{(2d) Fast and slow factors at $\xi\to0$.} The phase ties $y'-y$ to $O(\xi)$; concretely,
in the variables $y=z+\xi(w-z)$, $y'=z+\xi(w'-z)$ (the exact map that Step 3 uses), $\xi\to0$
at fixed $(w,w',z)$ sends both parent points onto the point $z$: $y\to z$, $y'\to z$. The limit
treats the factors of \eqref{eq:dsoftz} very differently:
\begin{itemize}
\item[] \emph{fast:}\quad $K(y,y',z)=\dfrac{K(w,w',z)}{\xi^{2}}$ blows up --- this is what generates the logarithm;
\item[] \emph{slow:}\quad $\KK\;\to\;\mathcal K(x-z)\!\cdot\!\mathcal K(x'-z)=K(x,x',z)$\quad and\quad $W(y,y')\to W(z,z)$ --- smooth, finite limits.
\end{itemize}

\textbf{(2e) Define the counterterm by freezing the slow factors.} The subtraction term is
obtained from \eqref{eq:dsoftz} by \emph{one} rule: keep the phase and the fast factor
untouched, and replace the slow factors by their collapse-point values,
\begin{align}
\KK\;W(y,y')\;\longrightarrow\;K(x,x',z)\;W(z,z)\,,
\end{align}
with an overall minus sign so that it subtracts. Explicitly, per $\xi$,
\begin{align}\label{eq:sub}
\frac{d^{3}N^{\rm sub}_{\xi}}{d^{3}k} = \frac{1}{(2\pi)^{3}}\, \frac{g^{2}}{2\pi^{2}}\, \frac{1}{k^{+}}\,\int_{x,y,x',y',z}\, e^{-ik\cdot (y'-y)\,\frac{1}{\xi}}\;\frac{1}{\pi}\,K(y,y',z)\,K(x,x',z)\;W(z,z)\,,
\end{align}
and, summed over $\xi$ with the same measure, kernel weight and prefactor as \eqref{eq:dsoft},
\begin{align}\label{eq:S}
S \;\equiv\; -\,\frac{g^{2}}{8\pi^{2}}\int_{\bar\xi}^{1}\frac{d\xi}{\xi^{2}}\;2N_{c}\,\frac{1-\xi}{\xi}\;\frac{d^{3}N^{\rm sub}_{\xi}}{d^{3}k}
\;=\;-\,\frac{1}{(2\pi)^{3}}\,\frac{g^{4}N_{c}}{8\pi^{5}}\,\frac{1}{k^{+}}
\int_{\bar\xi}^{1}\frac{d\xi}{\xi^{2}}\,\frac{1-\xi}{\xi}
\int_{x,y,x',y',z} e^{-ik(y'-y)\frac1\xi}\,K(y,y',z)\,K(x,x',z)\,\Wz\,.
\end{align}
Nothing new has been invented: the $\frac{1}{\pi}$, the prefactors and the measure of
\eqref{eq:sub}--\eqref{eq:S} are all inherited from \eqref{eq:dsoftz}. Physically,
\eqref{eq:sub} is the real-emission structure with the parent at the common point $z$ --- the
two soft-emission (dipole) kernels $K(y,y',z)\,K(x,x',z)$ of amplitude and conjugate, and the
Wilson lines at coincident points --- precisely the building blocks of the real BFKL kernel.
It is IR and UV finite in $d=2$ and contains no $\mu$.

\textbf{(2f) What the pair achieves.} Adding \eqref{eq:dsoftz} and \eqref{eq:S}, the two
integrands sit under the same integral sign and combine into a single bracket:
\begin{align}\label{eq:combo}
D_{\rm soft}+S
=+\,\frac{g^{2}}{8\pi^{2}}\int_{\bar\xi}^{1}\frac{d\xi}{\xi^{2}}\,2N_{c}\,\frac{1-\xi}{\xi}\,
\frac{1}{(2\pi)^{3}}\frac{g^{2}}{2\pi^{2}}\frac{1}{k^{+}}
\int_{x,y,x',y',z} e^{-ik\cdot(y'-y)\frac{1}{\xi}}\,\frac{K(y,y',z)}{\pi}
\Big[\KK\,W(y,y')-K(x,x',z)\,W(z,z)\Big]\,.
\end{align}
By (2d) the square bracket vanishes \emph{linearly} in $\xi$ at fixed $(w,w',z)$, and in those
variables the $z$-integral converges at large $z$ (the parent points travel with $z$, so both
terms of the bracket fall off together); so the pair is free of $\frac{2}{\epsilon}$,
$\gamma_{E}$ and $\mu$, and its soft region contributes
$\int\frac{d\xi}{\xi}\times O(\xi)={}$finite --- no $\log\frac{\vee}{k^{+}}$. The subtraction
$S$ alone therefore carries all the ultraviolet data and the entire rapidity logarithm of
$D_{\rm soft}$. And $S$, unlike $D_{\rm soft}$, can be computed in closed form: that is Step 3.

\section{Step 3: evaluating $S$ in closed form, explicitly}

\textbf{(3a) Only the phase knows about $\xi$.} Look at the inner integral of \eqref{eq:S} at
fixed $\xi$, with the $\xi$-independent factors pulled out of the $y,y'$ integration:
\begin{align}\label{eq:inner}
\frac{d^{3}N^{\rm sub}_{\xi}}{d^{3}k}
=\frac{1}{(2\pi)^{3}}\frac{g^{2}}{2\pi^{2}}\frac{1}{k^{+}}
\int_{x,x',z}\;K(x,x',z)\,W(z,z)\;\frac{1}{\pi}
\int_{y,y'} e^{-ik\cdot(y'-y)\frac{1}{\xi}}\;K(y,y',z)\,.
\end{align}
The variable $\xi$ appears in exactly one place: the phase. Strategy: change variables in the
$y,y'$ integral so that the phase loses its $\frac{1}{\xi}$; then every power of $\xi$ is
explicit, and the $\xi$-integral factorizes from the transverse one.

\textbf{(3b) The substitution.} At fixed $z$ set
\begin{align}\label{eq:rescale}
y=z+\xi\,(w-z)\;\Longleftrightarrow\;w=z+\frac{y-z}{\xi}\,,\qquad
y'=z+\xi\,(w'-z)\;\Longleftrightarrow\;w'=z+\frac{y'-z}{\xi}\,,
\end{align}
an exact, invertible, linear change of variables --- there is no approximation anywhere in
Step 3.

\textbf{(3c) Transform each factor, one at a time.}
\begin{itemize}
\item[(i)] Measure (two transverse dimensions each): $d^{2}y=\xi^{2}\,d^{2}w$,
$d^{2}y'=\xi^{2}\,d^{2}w'$, hence $d^{2}y\,d^{2}y'=\xi^{4}\,d^{2}w\,d^{2}w'$.
\item[(ii)] Differences: $y-z=\xi\,(w-z)$, $y'-z=\xi\,(w'-z)$, and, subtracting the first from
the second, $y'-y=\xi\,(w'-w)$.
\item[(iii)] Kernel:
\begin{align}
K(y,y',z)=\frac{(y-z)\cdot(y'-z)}{(y-z)^{2}\,(y'-z)^{2}}
=\frac{\xi^{2}\,(w-z)\cdot(w'-z)}{\xi^{2}(w-z)^{2}\;\xi^{2}(w'-z)^{2}}
=\frac{1}{\xi^{2}}\,K(w,w',z)\,.
\end{align}
\item[(iv)] Phase:
\begin{align}
e^{-ik\cdot(y'-y)\frac{1}{\xi}}=e^{-ik\cdot\xi(w'-w)\frac{1}{\xi}}=e^{-ik\cdot(w'-w)}\,.
\end{align}
The $\frac{1}{\xi}$ is eaten \emph{exactly} --- nothing is expanded or approximated.
\item[(v)] Spectators: $K(x,x',z)$, $W(z,z)$ and the $x,x',z$ measure contain no $y,y'$ and are
untouched.
\end{itemize}

\textbf{(3d) Assemble.} Putting (i)--(v) into \eqref{eq:inner},
\begin{align}\label{eq:nsubxi2}
\frac{d^{3}N^{\rm sub}_{\xi}}{d^{3}k}
=\xi^{4}\cdot\frac{1}{\xi^{2}}\cdot
\frac{1}{(2\pi)^{3}}\frac{g^{2}}{2\pi^{2}}\frac{1}{k^{+}}\,\frac{1}{\pi}\;\mathcal I
=\xi^{2}\;\frac{1}{(2\pi)^{3}}\frac{g^{2}}{2\pi^{2}}\frac{1}{k^{+}}\,\frac{1}{\pi}\;\mathcal I\,,
\qquad
\mathcal I\equiv\int_{x,x',z,w,w'} e^{-ik\cdot(w'-w)}\,K(w,w',z)\,K(x,x',z)\,W(z,z)\,,
\end{align}
and $\mathcal I$ is manifestly $\xi$-independent: \emph{all} the $\xi$-dependence of
$\frac{d^{3}N^{\rm sub}_{\xi}}{d^{3}k}$ is the explicit factor $\xi^{2}$ (Jacobian $\xi^{4}$
times kernel rescaling $\xi^{-2}$).

\textbf{(3e) The $\xi$-integral and the prefactor.} Insert \eqref{eq:nsubxi2} into
\eqref{eq:S}: the explicit $\frac{1}{\xi^{2}}$ of the measure is cancelled by the $\xi^{2}$ of
\eqref{eq:nsubxi2}, leaving
\begin{align}\label{eq:softint}
\int_{\bar\xi}^{1}d\xi\,\frac{1-\xi}{\xi}
=\int_{\bar\xi}^{1}\frac{d\xi}{\xi}-\int_{\bar\xi}^{1}d\xi
=\Big[\log\xi\Big]_{\bar\xi}^{1}-\Big[\xi\Big]_{\bar\xi}^{1}
=\big(0-\log\bar\xi\big)-\big(1-\bar\xi\big)
=\log\frac{\vee}{k^{+}}-\Big(1-\frac{k^{+}}{\vee}\Big)
\end{align}
(upper minus lower in both brackets; positive, as required for a positive integrand), and the
prefactor assembles as
\begin{align}
\frac{g^{2}}{8\pi^{2}}\cdot 2N_{c}\cdot\frac{1}{(2\pi)^{3}}\,\frac{g^{2}}{2\pi^{2}}\,\frac{1}{k^{+}}\cdot\frac{1}{\pi}
=\frac{1}{(2\pi)^{3}}\,\frac{g^{4}N_{c}}{8\pi^{5}}\,\frac{1}{k^{+}}\,.
\end{align}
Therefore
\begin{align}\label{eq:Svalue}
S=-\,\frac{1}{(2\pi)^{3}}\,\frac{g^{4}N_{c}}{8\pi^{5}}\,\frac{1}{k^{+}}\,
\Big[\log\frac{\vee}{k^{+}}-\Big(1-\frac{k^{+}}{\vee}\Big)\Big]
\int_{x,x',z,w,w'}e^{-ik\cdot(w'-w)}\,K(w,w',z)\,K(x,x',z)\,\Wz\,.
\end{align}

\section{Step 4: add and subtract, and the BFKL term}
Insert $0=S-S$ into \eqref{eq:leftover}, keeping $S$ in its integral form \eqref{eq:S} and
replacing $-S$ by minus the closed form \eqref{eq:Svalue}. Since \eqref{eq:Svalue} carries an
overall minus sign, $-S$ is \emph{plus}
$\frac{1}{(2\pi)^{3}}\frac{g^{4}N_{c}}{8\pi^{5}}\frac{1}{k^{+}}
\big[\log\frac{\vee}{k^{+}}-(1-\frac{k^{+}}{\vee})\big]\int KKW(z,z)$. Its
$\log\frac{\vee}{k^{+}}$ part is the \textbf{BFKL contribution of this row},
\begin{align}\label{eq:bfkl}
\boxed{\;\frac{d^{3}N^{\rm BFKL}_{\rm G\text{-}II,\,R\,IV}}{d^{3}k} = +\frac{1}{(2\pi)^{3}}\,\frac{g^{4}\,N_{c}}{8\pi^{5}}\,\frac{1}{k^{+}}\,{\color{blue}\log\Big(\frac{\vee}{k^{+}}\Big)}\int_{x,x',z,w,w'}e^{-ik\cdot (w'-w)}\,K(w,w',z)\,K(x,x',z)\,
\left[U^{ce'}(z) - U^{ce'}(x')\right]\left[U^{ce}(z) - U^{ce}(x)\right]\rho^{e'}(x')\rho^{e}(x)\;}
\end{align}
moved to the BFKL file (it appears there verbatim). Its coefficient is manifestly UV finite and
$\mu$-independent, as a BFKL kernel must be. The $-(1-\frac{k^{+}}{\vee})$ part of $-S$ is the
boundary term of the finite block.

\section{Result: the finite terms}
Collecting everything,
\begin{align}
\frac{d^{3}N_{\rm left}}{d^{3}k}
=\frac{d^{3}N^{\rm BFKL}}{d^{3}k}\Big|_{\eqref{eq:bfkl}}
+\frac{d^{3}N^{\rm finite}}{d^{3}k}\,,
\end{align}
with
\begin{align}\label{eq:finite}
    \frac{d^{3}N^{\rm finite}}{d^{3}k}  & =
   \underbrace{{\color{violet}-\,\frac{1}{(2\pi)^{3}}\,\frac{g^{4}N_{c}}{8\pi^{4}}\,\frac{1}{k^{+}}\,\log\Big(1-\frac{k^{+}}{\vee}\Big)\int_{x,y,x',y'}e^{-ik\cdot(y'-y)}\,\KK\;\Leps\;\Wbr}}_{\textstyle \rm endpoint\ (plus\text{-}prescription\ convention)\ term}
    \notag \\ & \underbrace{-\,\frac{1}{(2\pi)^{3}}\,\frac{g^{4}\,N_{c}}{8\pi^{5}}\, \frac{1}{k^{+}}\,\int_{k^{+}/\vee}^{1}\,\frac{d\xi}{\xi^{2}}\,\frac{(1-\xi)}{\xi}\,\int_{x,y,x',y',z}\, e^{-ik\cdot (y'-y)\,\frac{1}{\xi}}\,K(y,y',z)\,K(x,x',z)\,
\left[U^{ce'}(z) - U^{ce'}(x')\right]\left[U^{ce}(z) - U^{ce}(x)\right]\rho^{e'}(x')\,\rho^{e}(x)}_{\textstyle \rm soft\ subtraction\ (=S)} \notag \\ &
    \underbrace{-\frac{1}{(2\pi)^{3}}\,\frac{g^{4}\,N_{c}}{8\pi^{5}}\,\frac{1}{k^{+}}\,\Big(1-\frac{k^{+}}{\vee}\Big)\int_{x,x',z,w,w'}e^{-ik\cdot (w'-w)}\,K(w,w',z)\,K(x,x',z)\,
\left[U^{ce'}(z) - U^{ce'}(x')\right]\left[U^{ce}(z) - U^{ce}(x)\right]\rho^{e'}(x')\rho^{e}(x)}_{\textstyle \rm boundary} \notag \\
&  \underbrace{+\,\frac{1}{(2\pi)^{3}}\,\frac{g^{4}N_{c}}{8\pi^{4}}\, \frac{1}{k^{+}}\,\int_{\frac{k^{+}}{\vee}}^{1}\,d\xi\, \frac{(1-\xi)}{\xi}\, \int_{x,y,x',y'}\,e^{-ik\cdot (y'-y)\frac{1}{\xi}}\,\Bigg[\KK -2\,\frac{(x-y)\cdot (y'-y)\, (x'-y')\cdot (y'-y)}{(x-y)^{2}\,(x'-y')^{2}\,(y'-y)^{2}}\Bigg]\Wbr}_{\textstyle \rm azimuthal}\,.
\end{align}

\subsection*{Bookkeeping}
\begin{itemize}
\item[(i)] (soft subtraction $S$) $+$ (boundary) $+$ (BFKL entry \eqref{eq:bfkl}) $=0$
identically: the construction is an exact add-and-subtract of \eqref{eq:Svalue}, and all three
objects are UV finite.
\item[(ii)] The role of $S$ in the finite block is to pair with the soft channel
$2N_{c}\frac{1-\xi}{\xi}\otimes\Leps$ of the \emph{already extracted} DGLAP term: per $\xi$,
their sum is the $\mu$-free combination \eqref{eq:combo}, whose bracket is $O(\xi)$, so the
pair contains no $\log\frac{\vee}{k^{+}}\times\Leps$. The DGLAP term alone and $S$ alone each
contain that unphysical double structure; only their placement in different named blocks makes
it look otherwise.
\item[(iii)] The endpoint term (first entry of \eqref{eq:finite}) is the difference between the
two plus-prescription conventions. If $\mathcal P_{gg}$ is used with the standard $[0,1]$ plus
distribution, this entry is absorbed into the DGLAP term and disappears from
\eqref{eq:finite}; nothing else changes.
\item[(iv)] (BFKL \eqref{eq:bfkl}) $+$ \eqref{eq:finite} reproduces the left-over
\eqref{eq:leftover} identically --- Steps 2--4 only added zero.
\end{itemize}
\end{document}
